\documentclass[12pt]{article}
\usepackage{sbc-template}
\usepackage[utf8]{inputenc}
\usepackage[T1]{fontenc}
\IfFileExists{brazil.ldf}{\usepackage[brazil]{babel}}{\usepackage[english]{babel}}
\usepackage{graphicx,url,booktabs,tabularx,amsmath}
\sloppy
\title{Servidor HTTP/1.1 com Qualidade de Serviço em C com Pthreads: terceira entrega}
\author{Giovana Piazza\inst{1}, Luis Felipe Maccalli\inst{1}, Vitor Moreira\inst{1}}
\address{Universidade Federal do Pampa (Unipampa)\\
Fundamentos e Avaliação de Redes de Computadores -- 2026/2}
\begin{document}
\renewcommand{\refname}{Referências}
\maketitle
\begin{abstract}
This report presents the second version of a concurrent, persistent HTTP/1.1
server written in C using POSIX threads. It implements per-IP rate limits,
admission control, shared rates among connections and a live terminal panel.
The revision validates configuration and reports aggregate per-IP throughput.
Original packet captures are retained separately from new automated tests,
which include complete downloads of the images referenced by the HTML page.
Evaluation using Wireshark and IPTraf and experiments on physical networks
remain pending and are not reported as completed.
\end{abstract}
\begin{resumo}
Este relatório apresenta a segunda versão de um servidor HTTP/1.1 concorrente
e persistente em C com Pthreads. A implementação acrescenta taxas por IP,
controle de admissão, divisão da taxa entre conexões e painel em tempo real.
A revisão valida a configuração e corrige a medição agregada por IP. As
capturas originais são preservadas separadamente dos novos testes automatizados,
que incluem downloads completos das imagens referenciadas no HTML. A avaliação
com Wireshark e IPTraf e os experimentos em redes físicas permanecem pendentes,
sem atribuição de resultados que não foram obtidos.
\end{resumo}

\section{Introdução}
A terceira entrega acadêmica, prevista para 19/10/2026, corresponde ao relatório
e ao código da segunda versão do software. O enunciado exige servidor HTTP em
Linux, C e Pthreads, com clientes concorrentes, conexões persistentes e mecanismos
de qualidade de serviço (QoS)~\cite{spec}. A primeira versão constitui a base de
HTTP e concorrência; a segunda inclui controle de envio por IP e monitoramento.

O servidor aplica o ritmo de envio na camada de aplicação, acima do TCP. Isso
não substitui o controle de congestionamento ou de fluxo do transporte. A
implementação é um subconjunto didático de HTTP/1.1, com GET, HEAD e arquivos
estáticos, e não uma implementação completa das RFCs~\cite{http,sem}.

\section{Metodologia e organização}
O pacote mantém a primeira versão em V1 e a segunda em V2. A revisão de 19/10
foi feita sobre a V2 entregue pela equipe: inspeção do código e do relatório,
comparação com os requisitos, reprodução de cenários e correção dos problemas
identificados. Foram preservados o PCAP, o CSV e os registros originais em
\texttt{evidencias/v2}. As novas evidências ficam em
\texttt{evidencias/revisao-v2}, sem substituir os números históricos.

O relatório fornecido pela equipe registra acompanhamento por Trello e divisão
de responsabilidades. Esses dados de gestão não foram verificados pela execução
técnica da revisão. O cronograma restante prevê completar a avaliação com as
ferramentas exigidas e preparar a validação com o docente até 26/10, seguida dos
slides até 03/12, conforme o enunciado.

\section{Desenvolvimento e correções}
\subsection{Taxas, concorrência e admissão}
A execução segue o formato
\texttt{./server PORTA DIRETORIO VAZAO\_MAX\_KBPS [ARQUIVO\_TAXAS]}.
O arquivo auxiliar associa IPv4 e taxa inteira em kbps, usando a convenção
decimal: 1 kbps = 1000 bits/s. Endereços não cadastrados usam 1000 kbps, desde
que essa reserva caiba na capacidade informada.

Um mutex protege a tabela de IPs, suas conexões e a soma reservada. Se o IP já
estiver ativo, a conexão divide a reserva existente; para um IP sem conexões
ativas, a admissão exige $V_{reservada}+R_{IP}\le V_{max}$. A última conexão
libera a reserva. Uma recusa gera resposta 503. O limite de 128 conexões também
protege os recursos do processo.

A revisão rejeita taxas zero, negativas, inválidas, maiores que
1.000.000.000 kbps, valores fora da representação numérica, IPv4 inválido,
IP duplicado e campos excedentes. A capacidade passa pela mesma validação
numérica. Comentários com \texttt{\#} e linhas vazias são aceitos. Um arquivo
explicitamente informado que não pode ser aberto interrompe a inicialização;
quando apenas o arquivo padrão não existe, permanece a taxa padrão.

\subsection{Envio e divisão da taxa}
HTML é enviado sem limitação. Os demais objetos são enviados em blocos de até
8192 bytes. A taxa individual é recalculada a cada bloco por
$R_c=R_{IP}/N$, com divisão em ponto flutuante. Isso evita truncamento e a
imposição indevida de uma taxa mínima de 1 kbps para cada conexão.
Após enviar $n$ bytes, a espera nominal é
\begin{equation}
t_{espera}=\frac{8n}{1000R_c}\quad\text{segundos}.
\end{equation}
A espera usa \texttt{nanosleep}, fora do mutex e em fatias de até 100 ms,
permitindo encerramento mesmo com taxas muito baixas. O controle admite uma
rajada inicial de um bloco; o tempo de leitura e envio não é descontado da espera.
Portanto, não garante uma taxa instantânea rígida em todas as janelas de tempo.

\subsection{Painel e estimativas}
Uma thread atualiza o painel aproximadamente a cada segundo. Para cada IP,
mostra conexões, taxa reservada, taxa por conexão, estimativas e bytes de objetos
enviados durante o atendimento atual. A coluna REAL agora usa o total de bytes
de todas as conexões daquele IP na janela entre atualizações:
\begin{equation}
V_{IP}=\frac{8\,\Delta B_{IP}}{1000\,\Delta t}\quad\text{kbps},
\end{equation}
com $\Delta t$ em segundos. O contador é reiniciado quando um IP volta a ser
admitido após ficar sem conexões. A versão anterior mostrava a média de uma
transferência individual em uma linha que representava o IP completo.

O intervalo entre a solicitação do HTML e a do primeiro objeto do mesmo IP é
usado como estimativa didática de RTT, com resolução de milissegundos e mínimo
de 1 ms. O painel mostra também metade desse intervalo e o tamanho do HTML em
bits dividido pelo intervalo. Esses números incluem transmissão e processamento
do cliente: não medem diretamente RTT TCP, atraso unidirecional ou capacidade
do enlace. Navegações simultâneas do mesmo IP podem interferir na associação.

\section{Avaliação experimental}
\subsection{Evidências originais da equipe}
Os resultados abaixo são os registros fornecidos com a V2 original, anteriores
às correções. O relatório original informa execução em Zorin OS, GCC 13.3.0,
curl e tcpdump, na interface loopback. Essa configuração foi informada pela
equipe e não corresponde necessariamente ao ambiente da revisão.

\begin{table}[ht]
\centering\small
\caption{Resultados originais, objeto de 1 MiB por transferência.}
\begin{tabular}{llrr}
\toprule
Cenário & IP & Status & Vazão útil (kbps)\\
\midrule
Isolado & 127.0.0.1 & 200 & 2001,2\\
Três IPs & 127.0.0.1 & 200 & 2003,2\\
Três IPs & 127.0.0.3 & 200 & 1503,4\\
Três IPs & 127.0.0.4 (padrão) & 200 & 1004,7\\
Admissão & 127.0.0.2 & 503 & --\\
Mesmo IP, conexão 1 & 127.0.0.1 & 200 & 1004,3\\
Mesmo IP, conexão 2 & 127.0.0.1 & 200 & 1004,3\\
\bottomrule
\end{tabular}
\end{table}

A inspeção programática do PCAP original confirmou 1647 pacotes, 8 fluxos TCP,
8 SYN, 8 SYN-ACK, 14 requisições GET, 12 respostas 200 e 2 respostas 503,
compatíveis com o relatório fornecido. Essa inspeção não foi feita na interface
do Wireshark e não substitui a avaliação solicitada com essa ferramenta.
O script original usa um objeto binário não referenciado no HTML; por isso,
a revisão acrescenta o teste da página com as imagens efetivamente referenciadas.

\subsection{Testes da versão corrigida}
A revisão compilou o servidor com C11, otimização O2, Pthreads e opções
\texttt{-Wall -Wextra -Wpedantic -Werror}. Os testes abrem sockets TCP reais
no loopback. A suíte verifica configuração, inicialização, persistência,
integridade, admissão, liberação de reservas, taxa padrão, divisão entre
conexões e vazão agregada do painel. Os resultados detalhados e logs estão em
\texttt{evidencias/revisao-v2}.

Os 17 testes foram aprovados. A configuração inválida foi rejeitada antes da abertura do serviço. No teste com imagens, foram recebidos 50.331.864 bytes de conteúdo em 4,6505 s, com vazão útil agregada de 86,58 Mbps. A persistência e os hashes SHA-256 das quatro transferências foram confirmados.

No teste separado de QoS, a taxa padrão de 1000 kbps resultou em 936,2 kbps. Duas conexões do IP com reserva de 2000 kbps atingiram 950,4 e 944,1 kbps. As amostras agregadas do painel com duas conexões ficaram entre 1964,1 e 2097,2 kbps, em torno da reserva total de 2000 kbps. O teste de admissão observou 503 quando não havia reserva disponível e 200 após sua liberação.


O cenário de imagens usa dois clientes concorrentes do mesmo IP, com reserva
de 128000 kbps e capacidade de 128000 kbps. Cada cliente solicita HTML,
imagem1.bmp e imagem2.bmp na mesma conexão. Cada imagem possui 12.582.966 bytes,
verificados por SHA-256. A taxa elevada foi escolhida para testar arquivos
completos rapidamente; os testes separados de QoS usam taxas menores.

As medições são observações locais, sem caracterização de capacidade máxima ou
intervalos de confiança. A vazão útil do lote usa bytes de corpo recebidos e
tempo total até a conclusão dos dois clientes. Ela difere da vazão da interface,
que pode incluir cabeçalhos, retransmissões e tráfego de fundo.

\subsection{Pendências de ferramentas e redes físicas}
Wireshark, tcpdump e IPTraf não estão instalados no ambiente da revisão.
Não foram produzidas novas telas dessas ferramentas ou novas medições de
IPTraf. Nenhum resultado foi inventado para preencher essa ausência.
A avaliação exigida pelo enunciado permanece pendente até a execução no
computador da equipe. Ethernet e Wi-Fi com hosts distintos também não foram
avaliados; os vários endereços de loopback não representam hosts físicos.

O arquivo \texttt{V2/COMO-VALIDAR.md} e o script
\texttt{V2/avaliar-ferramentas.sh} orientam a captura com tcpdump, o uso do
Wireshark e a medição com IPTraf. O script verifica dependências e o processo
de captura, preservando os logs. A equipe deve registrar a interface, os hosts,
as unidades, a duração, o tráfego de fundo, retransmissões e, preferencialmente,
três repetições. Até isso ocorrer, não se afirma cumprimento integral da avaliação.

\section{Conclusões e próximos passos}
A versão corrigida implementa os mecanismos centrais de QoS e corrige falhas de
validação e de apresentação da vazão por IP. Os testes com arquivos completos
acrescentam evidência de persistência e integridade das imagens da página.
As taxas configuradas são reservas e alvos de envio de objetos; HTML, cabeçalhos
e erros não entram no limitador e não se afirma um teto físico absoluto.
Conexões ociosas participam da divisão enquanto permanecerem abertas.

Antes de submeter uma avaliação completa, a equipe deve executar e incorporar
as evidências reais de Wireshark/IPTraf. Na evolução, convém usar prazos de envio
cumulativos, melhorar a resolução temporal, separar sessões de navegação do mesmo
IP e repetir os cenários em redes físicas. Todos os integrantes devem compreender
as escolhas e conseguir demonstrar o código na validação com o docente.

\section{Apoio de IA e disponibilidade}
A V2 original declara apoio de Claude (Anthropic). A V1 original também registra
apoio de IA. As correções, os novos testes e esta documentação tiveram apoio de
OpenAI ChatGPT em 19/10/2026. Esses apoios estão registrados no código e README.
A equipe deve revisar as alterações e executar os experimentos complementares.
Não houve publicação automática no GitHub.

\begin{thebibliography}{9}
\bibitem{spec} Fundamentos e Avaliação de Redes de Computadores (2026).
Trabalho de programação em rede: especificação do trabalho prático. 2026/2.
\bibitem{forouzan} Forouzan, B. A.; Mosharraf, F. (2013). Redes de computadores:
uma abordagem top-down. AMGH.
\bibitem{http} Fielding, R.; Nottingham, M.; Reschke, J. (2022). HTTP/1.1.
RFC 9112. \url{https://www.rfc-editor.org/rfc/rfc9112}.
\bibitem{sem} Fielding, R.; Nottingham, M.; Reschke, J. (2022). HTTP Semantics.
RFC 9110. \url{https://www.rfc-editor.org/rfc/rfc9110}.
\bibitem{ia} OpenAI (2026). ChatGPT. Apoio à revisão em 19/10/2026.
\url{https://chatgpt.com/}. Anthropic (2026). Claude. Apoio declarado na V2 original.
\end{thebibliography}
\end{document}
